% ch10_e 第10章 §10.6尾–§10.7 (书页 492–504 / p507–p519)
% 实际覆盖范围（以 PNG 为准）：第10章正文止于书页 496（p511）末句
% "...and nothing more."；书页 497–501（p512–p516）为 BIBLIOGRAPHY，
% 书页 502–504（p517–p519）为 INDEX，均非正文，故本块正文实为书页 492–496。
% 块界说明：
%   起点——书页 492（p507）顶部 "strength of the U(1) gauge field." 为承接
%   书页 491 末句（ch10_d 末尾"而 $\pmb B$ 是"）的半句，故本块正文以
%   %JOIN% 起头续译。
%   终点——书页 496（p511）以 "...and nothing more." 自然收束第10章及全书
%   正文，其后为参考文献与索引，本块不设 TAIL 区。
% 本块无图、无习题；含脚注 1 条（原书脚注 78，书页 494）；带编号公式 2 条
% （(10.7.12)、(10.7.13)）。
%JOIN%
该 $U(1)$ 规范场的磁场强度。把它与 $U(1)$ 规范场论的标准哈密顿量，即 $H=g_1\pmb B^{2}+g_2\pmb E^{2}$（其中 $\pmb E$ 是 $U(1)$ 规范场的电场）相比较，我们发现我们的 $U(1)$ 哈密顿量缺少来自电场的动能项。动能项以我们模型中 $1/N_f$ 展开的高阶项出现。

为估计 $\pmb E^{2}$ 项中系数的值，我们注意到，对与时间无关的规范势，$\pmb E^{2}$ 具有 $(\partial_{\pmb x}a_0)^{2}$ 的形式。这样的项可以由式 (10.7.11) 积掉费米子而产生。产生的项量级为 $N_f\frac{1}{gN_f^{-1}|\chi|}(l_0\partial_{\pmb x}a_0)^{2}$。第一个系数 $N_f$ 来自如下事实：$N_f$ 族费米子中的每一族对 $(l_0\partial_{\pmb x}a_0)^{2}$ 的贡献都相同。第二个系数描写来自其中一族费米子的贡献。该系数具有能量倒数的量纲，其值可以用费米子带宽的倒数来估计，即 $1/(gN_f^{-1}|\chi|)$。于是，$U(1)$ 规范场的拉格朗日量具有如下形式（Marston and Affleck, 1989; Wen, 2002a）：
\begin{equation*}
\mathcal L=\frac{CN_f^{2}}{gl_0}\pmb E^{2}-(3l_0g|\chi|^{4})\pmb B^{2}
\end{equation*}
其中 $C$ 是一个 $O(1)$ 常数。这样的拉格朗日量描写我们 $SU(N_f)$ 自旋模型中的一种人工光。把上式与标准拉格朗日量 (10.6.18) 相比较，我们发现人工光的速度为 $c_a\sim l_0g/N_f$ 量级，而精细结构常数为 $1/N_f$ 量级。

在动量空间中，费米子跳跃哈密顿量 (10.7.9) 具有如下形式（注意对基态有 $a_0=0$）：
\begin{equation*}
H_f=\sum_{\pmb k}^{\prime}\Psi^{\dagger}_{a,\pmb k}\Gamma(\pmb k)\Psi_{a,\pmb k}
\end{equation*}
其中
\begin{gather*}
\Psi^{\top}_{a,\pmb k}=(\psi_{a,\pmb k},\psi_{a,\pmb k+\pmb Q_x},\psi_{a,\pmb k+\pmb Q_y},\psi_{a,\pmb k+\pmb Q_x+\pmb Q_y}),\\
\pmb Q_x=(\pi,0,0),\qquad \pmb Q_y=(0,\pi,0),\\
\Gamma(\pmb k)=-8|\chi|N_f^{-1}\big(\sin(k_x)\Gamma_1+\sin(k_y)\Gamma_2+\sin(k_z)\Gamma_3\big)\tag{10.7.12}
\end{gather*}
且 $\Gamma_1=\tau^{3}\otimes\tau^{0}$，$\Gamma_2=\tau^{1}\otimes\tau^{3}$，$\Gamma_3=\tau^{1}\otimes\tau^{1}$。这里 $\tau^{1,2,3}$ 是泡利矩阵，$\tau^{0}$ 是 $2\times2$ 单位矩阵。动量求和 $\sum_{\pmb k}^{\prime}$ 的取值范围为 $k_x\in(-\pi/2,\pi/2)$、$k_y\in(-\pi/2,\pi/2)$、$k_z\in(-\pi,\pi)$。由于 $\{\Gamma_i,\Gamma_j\}=2\delta_{ij}$（$i,j=1,2,3$），我们发现费米子具有色散
\begin{equation*}
E(\pmb k)=\pm8g|\chi|^{3}N_f^{-1}\sqrt{\sin^{2}(k_x)+\sin^{2}(k_y)+\sin^{2}(k_z)}
\end{equation*}
我们看到，该色散在 $\pmb k=0$ 与 $\pmb k=(0,0,\pi)$ 处有两个节点。因此，在连续极限下，式 (10.7.9) 将给出 $2N_f$ 个无质量的四分量狄拉克费米子。

在纳入 $U(1)$ 规范涨落之后，无质量的狄拉克费米子作为带单位电荷的费米子与 $U(1)$ 规范场相互作用。因此，我们 $SU(N_f)$ 自旋模型的总有效理论是一个具有 $2N_f$ 族单位电荷狄拉克费米子的 QED。我们将把这些费米子称为人工电子。连续有效理论具有如下形式：
\begin{equation*}
\mathcal L=\dot\psi_{I,a}D_0\gamma^{0}\psi_{I,a}+v_f\dot\psi_{I,a}D_i\gamma^{i}\psi_{I,a}+\frac{CN_f^{2}}{gl_0}\pmb E^{2}-(3l_0g|\chi|^{4})\pmb B^{2}+\dots\tag{10.7.13}
\end{equation*}
其中 $I=1,2$，$D_0=\partial_t+\mathrm{i}a_0$，$D_i=\partial_i+\mathrm{i}a_i|_{i=1,2,3}$，$v_f=8l_0g|\chi|^{3}/N_f^{-1}$，$\gamma_\mu|_{\mu=0,1,2,3}$ 是 $4\times4$ 狄拉克矩阵，且 $\bar\psi_{I,a}=\psi^{\dagger}_{I,a}\gamma^{0}$。这里 $\psi_{1,a}$ 与 $\psi_{2,a}$ 是构成 $SU(N_f)$ 基础表示的狄拉克费米子场。我们想指出，虽然人工光的速度 $c_a$ 与人工电子的速度 $v_f$ 都是 $l_0g/N_f$ 量级，但在我们的模型中这两个速度不必相同。因此，洛伦兹对称性并无保证。（译注：式 (10.7.13) 与 $v_f$ 均按原书排印转写；按此处 $\bar\psi_{I,a}$ 的定义，式中前两个 $\psi$ 上的点号疑为 $\bar\psi$ 横杠之误印，$v_f$ 表达式中的除号按色散关系疑为乘号之误，即 $v_f=8l_0g|\chi|^{3}N_f^{-1}$。）

式 (10.7.13) 描写 $SU(N_f)$ 模型在一个量子有序相——$\pi$ 通量相——中的低能动力学。费米子与规范玻色子都是无质量的，并且彼此相互作用。这里我们想讨论一个重要问题：在积掉高能的费米子涨落与规范涨落之后，费米子与规范玻色子还保持无质量吗？一般而言，无质量激发之间的相互作用会产生质量项，除非这种无质量性受到对称性或其他某种机制的保护。对我们的 $SU(N_f)$ 模型来说，基态不破坏任何对称性，所以我们不能用自发破缺的对称性来解释这些无质量激发。无质量激发受到刻画基态中量子序（或弦网凝聚）的 PSG 的保护（见 9.10 节及 Wen (2002a)）。

\subsection{评注：关于规范理论与费米统计的一些历史评注}

\begin{keypoints}
\item 看待规范场有两种方式：把它视为局域相位不变性的几何对象，或视为一个关联系统的集体模式。
\item “规范”的含义。
\item 规范场和费米场的存在并不意味着规范玻色子和费米子会作为低能准粒子出现。
\end{keypoints}

第一个系统的规范理论是 Maxwell 的电磁理论。虽然人们为了表示电场和磁场而引进了矢量势 $A_\mu$，但 $A_\mu$ 的含义并不清楚。

规范场的概念由 Weyl 于 1918 年引进，他还提出矢量势 $A_\mu$ 是一个规范场。Weyl 的想法受到 Einstein 引力理论的启发，是统一电磁学与引力的一次尝试。在 Einstein 的广义相对论中，坐标不变性导致了引力。于是 Weyl 想到，另一种几何对象的不变性或许会导致电磁学。他提出的是标度不变性。

考虑一个取值为 $f$ 的物理量。我们知道，若不指明单位，数值 $f$ 本身是没有意义的。用 $\omega$ 表示单位，这个物理量实际上由 $f\omega$ 给出。这就是标度的相对性。现在假设该物理量在空间每一点都有定义（这样我们考虑的就是一个物理场）。我们想知道，如何比较不同点 $x^{\mu}$ 与 $x^{\mu}+\mathrm{d}x^{\mu}$ 处的物理量。我们不能只比较数值 $f(x^{\mu})$ 与 $f(x^{\mu}+\mathrm{d}x^{\mu})$，因为单位 $\omega$ 在不同的点可能不同。对邻近的点 $x^{\mu}$ 与 $x^{\mu}+\mathrm{d}x^{\mu}$，两个单位仅相差一个接近 1 的因子，我们可以把这样的因子写成 $1+S_\mu\mathrm{d}x^{\mu}$。在 $x^{\mu}$ 与 $x^{\mu}+\mathrm{d}x^{\mu}$ 两处物理量之差不是由 $f(x^{\mu}+\mathrm{d}x^{\mu})-f(x^{\mu})=\partial_\mu f\,\mathrm{d}x^{\mu}$ 给出，而是由 $f(x^{\mu}+\mathrm{d}x^{\mu})(1+S_\mu\mathrm{d}x^{\mu})-f(x^{\mu})=(\partial_\mu+S_\mu)f\,\mathrm{d}x^{\mu}$ 给出。Weyl 证明，局域标度不变性要求只有 $S_\mu$ 的旋度才具有物理意义，正如在 Maxwell 理论中只有 $A_\mu$ 的旋度才有意义一样。这样，Weyl 把 $S_\mu$ 认作矢量势 $A_\mu$。Weyl 把局域标度不变性称为“Eich Invarianz”，这个词被译成了“规范不变性”（gauge invariance）。

然而，Weyl 的想法是错误的，矢量势 $A_\mu$ 并不能被认作那个“规范场” $S_\mu$。不过，Weyl 已经几乎对了。如果我们把物理场看作一个复波函数\footnote{复波函数的概念于 1925 年引进，比 Weyl 的“规范理论”晚了七年。}的振幅，把单位 $\omega$ 看作一个复相位，即 $|\omega|=1$，那么不同点处振幅之差便由 $(\partial_\mu+\mathrm{i}S_\mu)f\,\mathrm{d}x^{\mu}$ 给出，其中不同点的单位相差一个因子 $(1+\mathrm{i}S_\mu\mathrm{d}x^{\mu})$。正是这样的 $S_\mu$ 才能被认作矢量势。所以 $A_\mu$ 实际上应该叫作“相位场”，而“规范不变性”实际上应该叫作“相位不变性”。然而，旧名称一直沿用至今。

这段历史是试图赋予非物理的矢量 $A_\mu$ 某种物理（或几何）意义的一次尝试。它把矢量势视为纤维丛的联络。这一图像已被广泛接受。我们如今把矢量势称为规范场，把 Maxwell 理论称为规范理论。然而，这并不意味着我们必须从局域相位不变性出发把矢量势解释成一个几何对象。毕竟，量子波函数的相位是非物理的。

关于规范理论还有另一种观点。理论物理学界的许多思考者对规范势 $A_\mu$ 的冗余性颇为不满。早在 20 世纪 70 年代初人们就认识到，可以用规范不变的圈算符来刻画规范理论的不同相（Wegner, 1971; Wilson, 1974; Kogut and Susskind, 1975）。后来人们发现，可以用闭合弦来表述整个规范理论（Banks et al., 1977; Foerster, 1979; Gliozzi et al., 1979; Mandelstam, 1979; Polyakov, 1979; Savit, 1980）。这些研究揭示了规范理论与闭合弦理论之间的密切关系——这一观点与矢量势的几何观念大不相同。

在凝聚态物理的一项相关进展中，人们发现，如果玻色模型处于某些量子相，规范场可以从局域玻色模型中演生出来。这一现象也称为规范场的动力学产生。规范场从局域玻色模型中的演生有一段漫长而复杂的历史。演生的 $U(1)$ 规范场最早是在 $(1+1)$ 维 $CP^{N}$ 模型的量子无序相中引进的（D'Adda et al., 1978; Witten, 1979）。在凝聚态物理中，人们在正方格子上玻色自旋模型的自旋液体态的从属玻色子方法中发现了 $U(1)$ 规范场（Affleck and Marston, 1988; Baskaran and Anderson, 1988）。从属玻色子方法不仅含有 $U(1)$ 规范场，还含有无能隙的费米子场。然而，由于瞬子效应及其导致的 $U(1)$ 规范场在 $1+1$ 维和 $1+2$ 维中的禁闭（Polyakov, 1975），上述规范场与无能隙费米子场没有一个能给出作为低能物理准粒子出现的规范玻色子和无能隙费米子。即使在可以忽略瞬子效应的大 $N$ 极限下，$2+1$ 维中 $U(1)$ 规范场与无质量狄拉克费米子之间的边缘耦合也会破坏费米子传播子和规范传播子中的准粒子极点。这曾导致一种看法：$U(1)$ 规范场与无能隙费米子场只不过是“不可靠”的从属玻色子方法的一种非物理的人为产物。因此，寻找演生规范玻色子与演生费米子的关键，不在于写下一个含有规范场和费米场的拉格朗日量，而在于证明规范玻色子和费米子确实出现在物理的低能谱中。实际上，对任何给定的物理系统，我们总可以随意设计一个带有任意所选规范场的拉格朗日量来描写它。然而，拉格朗日量中的规范场未必给出一个作为低能准粒子出现的规范玻色子。只有当规范场的动力学使规范场处于解禁闭相时，规范玻色子才能作为低能准粒子出现。因此，在 D'Adda et al. (1978)、Witten (1979)、Baskaran and Anderson (1988) 与 Affleck and Marston (1988) 的最初发现之后，许多研究者一直在努力寻找规范场的解禁闭相。

在高能物理中，Foerster et al. (1980) 构造了一个具有演生的解禁闭 $U(1)$ 规范玻色子的 $(3+1)$ 维局域玻色模型。该工作提出，自然界中的光可能是演生的。在凝聚态物理中，有人证明，如果在一个二维自旋 1/2 模型中破坏时间反演对称性，那么由于 Chern--Simons 项的出现，从属玻色子方法中的 $U(1)$ 规范场可以处于解禁闭相（Khveshchenko and Wiegmann, 1989; Wen et al., 1989）。这一解禁闭相对应于自旋 1/2 模型的一个自旋液体态（Kalmeyer and Laughlin, 1987），称为手征自旋液体。第二种解禁闭相是通过把 $U(1)$ 规范结构破缺到 $Z_2$ 规范结构而找到的。这种相含有一个解禁闭的 $Z_2$ 规范理论（Read and Sachdev, 1991; Wen, 1991a），称为 $Z_2$ 自旋液体（或短程 RVB 态）。手征自旋液体和 $Z_2$ 自旋液体都有一些惊人的性质。准粒子激发携带自旋 1/2，对应于一次自旋翻转的\emph{一半}。尽管我们的自旋 1/2 模型是纯玻色模型，这些准粒子还可以携带分数统计或费米统计。这些凝聚态的例子表明，规范场和费米统计都可以从局域玻色模型中演生出来。

我们想指出，自旋液体并不是局域玻色模型演生费米子的第一个例子。演生费米子——或更一般地，演生任意子——的第一个例子由 FQH 态给出。虽然 Arovas et al. (1984) 只讨论了任意子如何从磁场中的费米子系统演生出来，但同样的论证很容易推广，用以表明费米子和任意子如何从磁场中的玻色子系统演生出来。此外，1987 年，在共振价键（RVB）态的研究中，人们在正方格子上的最近邻二聚体模型中提出了演生费米子（自旋子）（Kivelson et al., 1987; Rokhsar and Kivelson, 1988; Read and Chakraborty, 1989）。然而，按照解禁闭的图像，Kivelson et al. (1987) 与 Rokhsar and Kivelson (1988) 的结果只在二聚体模型的基态处于 $Z_2$ 解禁闭相时才成立。看来，只含最近邻二聚体的正方格子上的二聚体液体并不是解禁闭态（Rokhsar and Kivelson, 1988; Read and Chakraborty, 1989），因此，正方格子上的最近邻二聚体模型（Rokhsar and Kivelson, 1988）是否具有费米型准粒子仍不清楚（Read and Chakraborty, 1989）。不过，在三角格子上，二聚体液体确实是一个 $Z_2$ 解禁闭态（Moessner and Sondhi, 2001）。因此，Kivelson et al. (1987) 与 Rokhsar and Kivelson (1988) 的结果对三角格子二聚体模型成立，费米型准粒子确实在三角格子上的二聚体液体中演生出来。

上面所有具有演生费米子的模型都是 $(2+1)$ 维模型，在其中，演生费米子可以通过把通量束缚到带电粒子上来理解（Arovas et al., 1984）。最近，Levin and Wen (2003) 指出，演生费米子的关键是一种弦结构。在任意维数中，费米子一般都可以作为开弦的端点出现。这一弦图像使我们能够构造一个具有演生费米子的 $(3+1)$ 维局域玻色模型（Levin and Wen, 2003）。既然规范玻色子与费米子都可以作为弦网凝聚的结果而演生，我们可以说，弦网凝聚提供了统一规范玻色子与费米子的一条途径。

把二维正方格子上的玻色 $SU(N)$ 自旋模型（Affleck and Marston, 1988）加以推广，人们发现，从一个三维立方格子上的玻色 $SU(N)$ 自旋模型中，无能隙的解禁闭 $U(1)$ 规范玻色子和无能隙费米子都演生出来了（Wen, 2002a）。在 $1+3$ 维中，这两种无能隙激发可以彼此分离，因为它们在低能下相互作用很弱。$U(1)$ 规范玻色子和无能隙费米子在各个方面都表现得像光子和电子。因此，这个玻色 $SU(N)$ 自旋模型不仅包含人工光，还包含人工电子。

在规范理论和费米统计诞生约一百年之后，我们正面临如下问题。规范场的起源是什么——是几何的，还是动力学的？费米统计的起源是什么——是与生俱来的，还是演生的？在本书中，我们支持规范玻色子与费米子的动力学演生起源。规范玻色子和费米统计也许只不过是量子多玻色子系统的集体现象，仅此而已。
% ---- 第10章及全书正文至此结束（书页 496 末"...and nothing more."）。
% ---- 其后为 BIBLIOGRAPHY（书页 497–501）与 INDEX（书页 502–504），不属正文翻译范围。
